\documentclass[a4paper,11pt]{article}

% Packages for formatting and functionality
\usepackage[utf8]{inputenc}
\usepackage{geometry}
\usepackage{amsmath, amssymb, amsfonts, amsthm}
\usepackage{graphicx}
\usepackage{hyperref}
\usepackage{authblk}
\usepackage{booktabs}
\usepackage{titlesec}
\usepackage{caption}
\usepackage{subcaption}
\usepackage{float}
\usepackage{xcolor}
\usepackage{algorithm}
\usepackage{algpseudocode}
\graphicspath{{assets/}}

% Theorem environments
\newtheorem{theorem}{Theorem}
\newtheorem{proposition}{Proposition}
\newtheorem{lemma}{Lemma}
\newtheorem{corollary}{Corollary}
\newtheorem{definition}{Definition}
\newtheorem{remark}{Remark}

% Margin settings
\geometry{
    top=1in,
    bottom=1in,
    left=1in,
    right=1in
}

% Title and Author Info
\title{\textbf{Holographic Invariant Storage: Design-Time Safety Contracts\\via Vector Symbolic Architectures}}
\author{Arsenios Scrivens}
\affil{Independent Researcher}
\date{March 13, 2026}

\begin{document}

\maketitle

\begin{abstract}
We introduce Holographic Invariant Storage (HIS), a protocol that assembles known properties of bipolar Vector Symbolic Architectures into a \textit{design-time safety contract} for LLM context-drift mitigation. The contract provides three closed-form guarantees evaluable before deployment: single-signal recovery fidelity converging to $1/\sqrt{2} \approx 0.707$ (regardless of noise depth or content), continuous-noise robustness $2\Phi(1/\sigma) - 1$, and multi-signal capacity degradation $\approx\!\sqrt{1/(K{+}1)}$. These bounds, validated by Monte Carlo simulation ($n = 1{,}000$), enable a systems engineer to budget recovery fidelity and codebook capacity at design time---a property no timer or embedding-distance metric provides. A pilot behavioral experiment (four LLMs, 2B--7B, 720 trials) confirms that safety re-injection improves adherence at the 2B scale; full results are in Appendix~\ref{sec:llm_experiment}.
\end{abstract}

\textbf{Keywords:} Vector Symbolic Architectures, Hyperdimensional Computing, LLM safety, context drift, design-time contract, bipolar hypervectors

\section{Introduction}

Context drift---the degradation of adherence to an original system prompt as the context window fills with interaction history---is a \textit{conditional} problem. It binds when a model's intrinsic alignment is insufficient for its deployment context: smaller or fine-tuned models, longer sessions, adversarial prompt sequences, or safety-critical domains where even marginal degradation is unacceptable. Under standard conversational conditions, frontier-scale models (GPT-4-class and beyond, as of March 2026) maintain robust safety adherence, and our own experiments show ceiling-level safety at 7B parameters (Appendix~\ref{sec:llm_experiment}). External mitigation mechanisms are therefore relevant primarily in the gap between a model's intrinsic alignment capability and the demands of its deployment context.

For deployments where that gap exists, the underlying vulnerability is architectural: the attention mechanism treats safety constraints as just another token sequence, weighted probabilistically against the immediate conversational context \cite{vaswani2017, liu2024}. Positional recency bias in causal attention can cause early instructions---including safety constraints---to be ``lost in the middle'' of long sequences, and modern context windows (128K--1M tokens) increase the volume of competing signals. This makes safety adherence load-dependent: it varies with context length, noise content, and prompt distribution \cite{wei2023, perez2022}.

Current approaches to this problem fall into several categories. Reinforcement Learning from Human Feedback (RLHF) \cite{ouyang2022} and Constitutional AI \cite{bai2022} embed safety preferences into model weights during training, providing baseline robustness but offering no mechanism for runtime verification or recovery when context-level drift occurs. Retrieval-Augmented Generation (RAG) \cite{lewis2020} and scratchpad-based memory \cite{nye2021} provide external memory, but their retrieved content still competes with context noise through the attention mechanism. Periodic re-prompting (re-injecting the system prompt at intervals) is a common engineering heuristic but lacks theoretical grounding and scales poorly with prompt length.

We propose a complementary approach: storing the safety signal in an external memory substrate based on \textbf{Vector Symbolic Architectures (VSA)} \cite{kanerva2009, plate1995, gayler2003}. This paper is a \textit{theoretical framework contribution}: we assemble known properties of bipolar hyperdimensional computing into a \textit{design-time safety contract}---a set of closed-form guarantees that a systems engineer can evaluate before deployment---and validate it at the signal level via Monte Carlo simulation. A pilot behavioral experiment is included as an appendix to demonstrate end-to-end feasibility but is not the paper's primary contribution.

HIS addresses one specific sub-problem: \textit{signal storage and recovery}---ensuring that a safety instruction can be retrieved with known fidelity from a corrupted representation. It does \textbf{not} solve the downstream \textit{signal utilization} problem: the recovered instruction is re-injected through the standard context window, where it competes for attention weight like any other token sequence. This scope limitation applies equally to all re-injection-based approaches (periodic re-prompting, RAG retrieval, embedding-triggered re-injection); HIS's contribution is providing algebraic guarantees on the \textit{storage} side that these alternatives lack.

Hyperdimensional Computing (HDC) provides algebraic operations over high-dimensional distributed representations that exhibit well-characterized noise tolerance properties \cite{kanerva2009, thomas2021}. In $D$-dimensional bipolar spaces, randomly generated vectors are near-orthogonal with high probability \cite{kanerva2009, vershynin2018}, meaning that a safety signal encoded as a hypervector can be recovered from additive noise via algebraic inversion rather than statistical inference. The VSA pipeline additionally provides a \textit{continuous drift metric}: the pre-restoration cosine similarity between the current context encoding and the stored invariant degrades proportionally to context corruption, signalling \textit{when} re-injection is needed---information that a static text file or fixed-interval timer cannot provide.

\subsection{Contributions}
\begin{enumerate}
    \item \textbf{Protocol Specification:} We present Holographic Invariant Storage (HIS), a VSA-based protocol for encoding, storing, and restoring safety constraints with a continuous drift metric (Section~3).
    \item \textbf{Design-Time Safety Contract:} We assemble standard properties of bipolar superposition \cite{kanerva2009} into closed-form fidelity guarantees evaluable \textit{before deployment}: single-signal recovery converging to $0.707$ (Theorem~\ref{thm:geometric_bound}), continuous-noise robustness (Proposition~\ref{prop:continuous_noise}), and multi-signal capacity scaling as $\approx\!\sqrt{1/(K{+}1)}$ (Proposition~\ref{prop:multi_signal}). While each bound is individually derivable from the HDC literature, no existing work synthesizes them into a deployable LLM safety contract or characterizes the contract's limits (normalization requirement, shared-variance ceiling, codebook threshold).
    \item \textbf{Signal-Level Validation:} Monte Carlo validation ($n = 1{,}000$) confirming convergence to theoretical bounds, signal-level comparisons against re-prompting, RAG, and no-intervention baselines (Appendix~\ref{sec:baselines}), and a 50-turn integration proof-of-concept with 100\% codebook retrieval (Appendix~\ref{sec:integration_poc}).
    \item \textbf{Pilot Behavioral Experiment:} As a feasibility demonstration, we deploy the full HIS pipeline with four open-source LLMs (2B--7B, 720 total trials). Re-injection improves safety adherence over no intervention at 2B ($p_{\text{uncorr}} = 0.013$, $p_{\text{Holm}} = 0.065$; $d = 0.71$). This pilot scopes effect sizes and conditions for a definitive behavioral study (Appendix~\ref{sec:llm_experiment}).
\end{enumerate}

Our contribution is the \textit{contract itself}---the guarantees, their limits, and the engineering framework that makes them evaluable at design time. The appendices provide supplementary material---extended proofs, robustness tests, signal-level baselines, and a pilot behavioral experiment---that support but are not the core of the contribution.

\section{Related Work}
\label{sec:related_work}

\textbf{AI Safety and Alignment.}
The concrete challenges of AI safety were catalogued by Amodei et al.\ \cite{amodei2016}, including reward hacking, distributional shift, and safe exploration. RLHF \cite{ouyang2022} addresses alignment at training time by shaping model weights to reflect human preferences, while Constitutional AI \cite{bai2022} automates this process via self-critique. These methods embed safety into model parameters but provide no runtime mechanism to detect or correct context-level drift during long inference sessions. Wei et al.\ \cite{wei2023} demonstrate that safety training can be circumvented via ``jailbreak'' attacks that exploit in-context reasoning---precisely the failure mode HIS is designed to complement.

\textbf{Representation Engineering and Activation Steering.}
A complementary line of work manipulates model internals directly to steer behavior. Turner et al.\ \cite{turner2023} introduce activation addition, injecting steering vectors into residual stream activations to control model outputs without fine-tuning. Zou et al.\ \cite{zou2023repeng} formalize this as representation engineering, identifying linear directions in activation space that correspond to high-level concepts (honesty, safety, harmfulness) and using them to read or control model behavior at inference time. These approaches operate at the representation level \textit{inside} the model, whereas HIS operates externally on the context window; the two are complementary. HIS could potentially be combined with activation steering by using the drift metric to trigger not only context-level re-injection but also activation-level intervention.

\textbf{External Memory and Context Management.}
Retrieval-Augmented Generation (RAG) \cite{lewis2020} augments LLMs with retrieved documents, but retrieved content enters the context window and is subject to the same attention-based dilution as any other token sequence. Scratchpad methods \cite{nye2021, wei2022cot} use the context window itself as working memory, offering no protection against context noise. Memory-augmented architectures \cite{graves2014, wu2022} provide differentiable external memory but are trained end-to-end and do not offer algebraic recovery guarantees. Recent work on KV-cache compression and attention-sink methods (e.g., StreamingLLM \cite{xiao2024}) addresses the ``lost in the middle'' problem at the attention level; these are complementary to HIS, which operates on the representation level outside the model.

\textbf{Vector Symbolic Architectures and Hyperdimensional Computing.}
VSAs were introduced by Plate \cite{plate1995} as Holographic Reduced Representations (HRR) and independently developed by Kanerva \cite{kanerva2009} and Gayler \cite{gayler2003}. These architectures exploit the concentration of measure in high-dimensional spaces \cite{vershynin2018}: in $D \geq 1{,}000$ dimensions, random bipolar vectors are near-orthogonal ($|\cos\theta| \approx 1/\sqrt{D}$), enabling robust associative memory retrieval. Recent work has applied HDC to classification \cite{thomas2021}, language processing \cite{schlegel2022}, and cognitive modeling \cite{eliasmith2012}. The noise tolerance properties we exploit---including the $1/\sqrt{2}$ recovery bound for bipolar superposition and the capacity scaling of bundled associative memories---are well-characterized in this literature. Our contribution is not the derivation of these properties but their \textit{synthesis into a deployable safety contract}: the specific identification of which bounds matter for LLM safety-signal preservation, the characterization of the contract's limits (normalization requirement, shared-variance ceiling, codebook-size threshold), and the worked engineering framework that makes these guarantees evaluable before deployment. This combination---applied to a domain (runtime LLM safety) that the existing VSA literature does not address---constitutes the paper's theoretical contribution.

\textbf{Guardrails and Output Filters.}
Engineering frameworks such as NeMo Guardrails \cite{rebedea2023} provide rule-based or LLM-mediated output filtering to enforce safety constraints at inference time. These systems operate on the \textit{output} side (intercepting and filtering generated text), whereas HIS operates on the \textit{input} side (re-injecting safety instructions into the context). The approaches are complementary: guardrails catch unsafe outputs; HIS aims to prevent drift that produces them.

The idea of separating invariant memory from working memory also has precedent in cognitive architectures (SOAR \cite{laird2012}, ACT-R \cite{anderson2004}); HIS can be viewed as a minimalist instantiation of this principle.

\section{Methodology}

\subsection{Vector Symbolic Architecture (VSA)}
Our approach utilizes $D$-dimensional bipolar hypervectors ($v \in \{-1,1\}^{D}$, with $D = 10{,}000$ throughout) to represent semantic concepts. We rely on three algebraic operations whose properties are well-established in the VSA literature \cite{kanerva2009, plate1995, gayler2003}:

\begin{itemize}
    \item \textbf{Binding ($\otimes$):} Element-wise multiplication of two vectors. For bipolar vectors, $a \otimes b \in \{-1,1\}^D$, and the result is near-orthogonal to both inputs ($\mathbb{E}[\cos(a, a \otimes b)] \approx 0$). Binding is its own inverse: $a \otimes a = \mathbf{1}$, so $(a \otimes b) \otimes a = b$. This is the mechanism that enables key-based retrieval.
    \item \textbf{Bundling ($+$):} Element-wise addition, creating a superposition that is similar to each of its components. This represents the accumulation of context signals (including noise).
    \item \textbf{Cleanup ($\text{sign}(\cdot)$):} Element-wise sign function applied to a real-valued superposition, following the standard numerical convention:
    \begin{equation}
        \text{sign}(x) = \begin{cases} +1 & \text{if } x > 0 \\ 0 & \text{if } x = 0 \\ -1 & \text{if } x < 0 \end{cases}
        \label{eq:sign_convention}
    \end{equation}
    This convention is implemented by all standard numerical libraries (NumPy, PyTorch, Julia). The $\text{sign}$ function serves as the primary noise-suppression mechanism: dimensions where signal and noise agree are retained, while dimensions where they cancel are \textit{abstained} ($\text{sign}(0) = 0$). The consequences of this convention for recovery fidelity are analyzed in Section~4 (Remark~\ref{rem:sign_convention}).
\end{itemize}

\subsection{The Restoration Protocol}
We define the agent's safety constraint as a ``System Invariant'' ($H_{\text{inv}}$), created by binding a known Goal Key ($K_{\text{goal}}$) to the Safe Value vector ($V_{\text{safe}}$):
\begin{equation}
    H_{\text{inv}} = K_{\text{goal}} \otimes V_{\text{safe}}
    \label{eq:invariant}
\end{equation}

During operation, this invariant is corrupted by additive context noise ($N_{\text{context}}$), producing a drifted state. The restoration protocol recovers the original value by: (1) normalizing and superimposing the noise, (2) applying binarization, and (3) unbinding with the original key:
\begin{equation}
    V_{\text{recovered}} = \text{sign}\!\left(H_{\text{inv}} + \hat{N}_{\text{context}}\right) \otimes K_{\text{goal}}
    \label{eq:restoration}
\end{equation}
where $\hat{N}_{\text{context}}$ denotes the noise after normalization (see Section~\ref{sec:normalization}).

The recovery works because $K_{\text{goal}}$ is near-orthogonal to $\hat{N}_{\text{context}}$ in high dimensions, so unbinding distributes the noise uniformly across the hyperspace (each component contributes $\approx 0$ in expectation), while the signal component reconstructs $V_{\text{safe}}$ coherently.

\begin{algorithm}[H]
\caption{Holographic Invariant Storage: Restoration Protocol}\label{alg:restore}
\begin{algorithmic}[1]
\State \textbf{Offline (Initialization):}
\State Generate random bipolar key: $K_{\text{goal}} \in \{-1,1\}^D$
\State Encode safety constraint: $V_{\text{safe}} \in \{-1,1\}^D$
\State Compute invariant: $H_{\text{inv}} \leftarrow K_{\text{goal}} \otimes V_{\text{safe}}$
\State Store $(K_{\text{goal}}, H_{\text{inv}})$ in external memory (outside context window)
\Statex
\State \textbf{Online (Each Restoration Step):}
\State Encode current context as noise vector: $N_{\text{context}} \in \mathbb{R}^D$
\State Normalize: $\hat{N}_{\text{context}} \leftarrow N_{\text{context}} \cdot \frac{\|H_{\text{inv}}\|}{\|N_{\text{context}}\|}$ \Comment{See Section~\ref{sec:normalization}}
\State Superimpose: $S \leftarrow H_{\text{inv}} + \hat{N}_{\text{context}}$
\State Binarize: $S_{\text{clean}} \leftarrow \text{sign}(S)$
\State Unbind: $V_{\text{recovered}} \leftarrow S_{\text{clean}} \otimes K_{\text{goal}}$
\State \textbf{Return} $V_{\text{recovered}}$ \Comment{Approximate recovery of $V_{\text{safe}}$}
\end{algorithmic}
\end{algorithm}

\subsection{The Normalization Constraint}
\label{sec:normalization}
The geometric bound (Theorem~\ref{thm:geometric_bound}) holds under the condition that the noise vector is normalized to have the same magnitude as the invariant: $\|\hat{N}_{\text{context}}\| = \|H_{\text{inv}}\|$. For bipolar vectors, $\|H_{\text{inv}}\| = \sqrt{D}$, so this requires scaling the noise to $\|\hat{N}_{\text{context}}\| = \sqrt{D}$.

This is a \textbf{design constraint}, not a natural property of the system. In practice, context noise accumulates over time, and its raw magnitude may grow unboundedly. The normalization step (Algorithm~\ref{alg:restore}, Line 8) enforces this constraint explicitly. This is analogous to gain control in signal processing or batch normalization in neural networks \cite{ioffe2015}---a deliberate architectural choice that keeps the signal-to-noise ratio fixed at 1:1 (0 dB).

This design choice has consequences:
\begin{itemize}
    \item \textbf{Advantage:} Recovery fidelity becomes a deterministic geometric property of the architecture, independent of noise magnitude or content.
    \item \textbf{Limitation:} The normalization discards information about the \textit{magnitude} of the noise. A system under heavy attack and a system under mild perturbation produce the same normalized noise vector. The protocol recovers the safety signal with equal fidelity in both cases---but it does not provide a \textit{measure of threat level}. A complementary anomaly detection mechanism would be needed for that purpose.
\end{itemize}

\section{Theoretical Analysis}

We now characterize the expected recovery fidelity of the restoration protocol. The result below is a direct consequence of well-known properties of bipolar superposition in HDC \cite{kanerva2009}; we state it as a theorem for self-containedness and to fix notation for the extensions that follow (Propositions~\ref{prop:continuous_noise}--\ref{prop:multi_signal}), not as a claim of mathematical novelty. The relevant property is that the standard $\text{sign}$ convention---$\text{sign}(0) = 0$ (Equation~\ref{eq:sign_convention})---introduces \textit{abstention} on tied dimensions, which reduces the norm of the cleaned vector and produces a cosine similarity strictly above what a random-tiebreaking convention would yield.

\begin{theorem}[Geometric Recovery Bound]
\label{thm:geometric_bound}
Let $H_{\text{inv}} \in \{-1,1\}^D$ and $\hat{N}_{\text{context}} \in \{-1,1\}^D$ be independent, uniformly random bipolar vectors with $D \gg 1$. Define $S = H_{\text{inv}} + \hat{N}_{\text{context}}$ and $S_{\text{clean}} = \text{sign}(S)$ under the standard convention (Equation~\ref{eq:sign_convention}). Then:
\begin{equation}
    \mathbb{E}\!\left[\text{CosSim}(S_{\text{clean}},\, H_{\text{inv}})\right] = \frac{1}{\sqrt{2}} \approx 0.7071
    \label{eq:geometric_bound}
\end{equation}
with concentration: let $K = |\{i : S_{\text{clean},i} \neq 0\}|$ be the number of agreement dimensions. Since $\text{CosSim}(S_{\text{clean}},\, H_{\text{inv}}) = \sqrt{K/D}$ (see proof), Hoeffding's inequality on $K/D$ gives
\begin{equation}
    P\!\left(\left|\frac{K}{D} - \frac{1}{2}\right| > t\right) \leq 2\exp\!\left(-2Dt^2\right)
    \label{eq:concentration_K}
\end{equation}
and Lipschitz transfer ($\sqrt{\cdot}$ has Lipschitz constant $L \leq 1$ on $[1/4,\,1]$) yields
\begin{equation}
    P\!\left(\left|\text{CosSim} - \frac{1}{\sqrt{2}}\right| > t\right) \leq 2\exp\!\left(-2Dt^2\right) \quad \text{for } 0 < t \leq \tfrac{1}{2\sqrt{2}}
    \label{eq:concentration}
\end{equation}
\end{theorem}

\begin{proof}
We analyze the per-dimension behavior, compute the inner product and norms, and assemble the cosine similarity.

\textbf{Step 1: Per-dimension outcomes.}
For each dimension $i$, the sum $S_i = H_{\text{inv},i} + \hat{N}_{\text{context},i}$ takes values in $\{-2,\, 0,\, +2\}$:
\begin{itemize}
    \item \textbf{Agreement} ($S_i = \pm 2$, probability $1/2$): Both vectors share the same sign. Then $\text{sign}(S_i) = H_{\text{inv},i}$ with certainty.
    \item \textbf{Cancellation} ($S_i = 0$, probability $1/2$): The two vectors disagree. Then $\text{sign}(S_i) = 0$ by the standard convention.
\end{itemize}
Thus $S_{\text{clean}}$ is a \textit{ternary} vector in $\{-1, 0, +1\}^D$, with each component independently equal to $H_{\text{inv},i}$ (with probability $1/2$) or $0$ (with probability $1/2$).

\textbf{Step 2: Inner product.}
The per-dimension contribution to the inner product is:
\begin{equation}
    S_{\text{clean},i} \cdot H_{\text{inv},i} = \begin{cases} H_{\text{inv},i}^2 = 1 & \text{with probability } 1/2 \\ 0 \cdot H_{\text{inv},i} = 0 & \text{with probability } 1/2 \end{cases}
\end{equation}
Therefore $\mathbb{E}[S_{\text{clean},i} \cdot H_{\text{inv},i}] = 1/2$, and by linearity:
\begin{equation}
    \mathbb{E}\!\left[\langle S_{\text{clean}},\, H_{\text{inv}} \rangle\right] = \frac{D}{2}
    \label{eq:inner_product}
\end{equation}

\textbf{Step 3: Norms.}
The norm of $H_{\text{inv}}$ is deterministic: $\|H_{\text{inv}}\| = \sqrt{D}$.

The norm of $S_{\text{clean}}$ depends on how many dimensions are non-zero. Let $K = |\{i : S_{\text{clean},i} \neq 0\}|$ be the number of agreement dimensions. Then $K \sim \text{Binomial}(D,\, 1/2)$ and:
\begin{equation}
    \|S_{\text{clean}}\|^2 = \sum_{i=1}^{D} S_{\text{clean},i}^2 = K
    \label{eq:norm_sclean}
\end{equation}
since each non-zero entry is $\pm 1$. Thus $\|S_{\text{clean}}\| = \sqrt{K}$, and by the law of large numbers, $K/D \to 1/2$ almost surely, giving $\|S_{\text{clean}}\| \to \sqrt{D/2}$.

\textbf{Step 4: Cosine similarity.}
Assembling the pieces:
\begin{align}
    \text{CosSim}(S_{\text{clean}},\, H_{\text{inv}}) &= \frac{\langle S_{\text{clean}},\, H_{\text{inv}} \rangle}{\|S_{\text{clean}}\| \cdot \|H_{\text{inv}}\|} \notag \\
    &= \frac{K}{\sqrt{K} \cdot \sqrt{D}} = \frac{\sqrt{K}}{\sqrt{D}} = \sqrt{\frac{K}{D}}
    \label{eq:cosine_derivation}
\end{align}
where we used the fact that $\langle S_{\text{clean}},\, H_{\text{inv}} \rangle = K$ (every non-zero entry of $S_{\text{clean}}$ agrees with $H_{\text{inv}}$ by construction, contributing $+1$; every zero entry contributes $0$). Taking expectations:
\begin{equation}
    \mathbb{E}\!\left[\text{CosSim}\right] = \mathbb{E}\!\left[\sqrt{K/D}\right] \to \sqrt{1/2} = \frac{1}{\sqrt{2}} \approx 0.7071
    \label{eq:final_bound}
\end{equation}
The convergence is justified by the concentration of $K/D$ around $1/2$: since $K \sim \text{Binomial}(D,\, 1/2)$, Hoeffding's inequality \cite{hoeffding1963} gives $P(|K/D - 1/2| > t) \leq 2\exp(-2Dt^2)$. The transfer to $\sqrt{K/D}$ follows because $f(x) = \sqrt{x}$ is Lipschitz with constant $L = 1/(2\sqrt{\epsilon})$ on $[\epsilon, 1]$; taking $\epsilon = 1/4$ (the high-probability region $K/D \geq 1/4$), we obtain $P(|\sqrt{K/D} - 1/\sqrt{2}| > t) \leq 2\exp(-2D t^2)$ for $t \leq 1/(2\sqrt{2})$, confirming sub-Gaussian concentration of the cosine similarity. (By Jensen's inequality for concave $\sqrt{\cdot}$, $\mathbb{E}[\sqrt{K/D}] \leq \sqrt{1/2}$; the gap is $O(1/D)$ and negligible for $D = 10{,}000$.)
\end{proof}

\begin{remark}[Why $\text{sign}(0) = 0$ Matters]
\label{rem:sign_convention}
The bound $1/\sqrt{2}$ depends critically on the sign convention. If ties were broken randomly (i.e., $\text{sign}(0) = \pm 1$ with equal probability), then $S_{\text{clean}} \in \{-1,+1\}^D$ and $\|S_{\text{clean}}\| = \sqrt{D}$, but the inner product remains $D/2$, yielding:
\[
    \text{CosSim}_{\text{random-tiebreak}} = \frac{D/2}{\sqrt{D} \cdot \sqrt{D}} = \frac{1}{2}
\]
The standard convention ($\text{sign}(0) = 0$) produces strictly higher fidelity ($0.707$ vs.\ $0.500$) because abstaining on uncertain dimensions \textit{reduces the denominator} without affecting the numerator---the zeros contribute nothing to either the inner product or the noise. This is not an implementation artifact; it is a fundamental property of the recovery geometry.
\end{remark}

\begin{remark}[Interpretation]
The $1/\sqrt{2}$ bound is a \textit{structural property of the architecture}, independent of noise content. It arises from a simple combinatorial fact: when two independent bipolar vectors are summed, exactly half their dimensions agree (in expectation). The $\text{sign}$ operation retains only the agreeing dimensions and discards the rest, producing a sparse vector that is geometrically closer to the signal than a dense noisy vector would be. The Monte Carlo simulation (Section~\ref{sec:empirical}) does not ``discover'' this bound; it \textit{validates} that the implementation correctly realizes it.
\end{remark}

\begin{remark}[Practical Meaning of 0.71 Fidelity]
\label{rem:practical_fidelity}
A cosine similarity of $0.71$ corresponds to ${\sim}50\%$ shared variance ($\cos^2(\pi/4) = 0.5$)---not ``high fidelity'' in isolation. However, for codebook retrieval in $D = 10{,}000$ dimensions, $0.71$ suffices to discriminate the correct vector from $K > 10^6$ alternatives: random bipolar vectors have pairwise similarity $\approx 0 \pm 0.01$ \cite{kanerva2009}, placing the recovered signal ${\sim}70$ standard deviations above the noise floor.
\end{remark}

\subsection{Extension to Continuous Noise}
\label{sec:continuous_noise}
Theorem~\ref{thm:geometric_bound} assumes bipolar noise ($\hat{N} \in \{-1,1\}^D$), which produces exact cancellation ($S_i = 0$) with probability $1/2$. In practice, the normalized noise vector produced by a semantic encoder is continuous-valued, not bipolar. We now characterize recovery under Gaussian continuous noise.

\begin{proposition}[Continuous Noise Recovery]
\label{prop:continuous_noise}
Let $H_{\text{inv}} \in \{-1,1\}^D$ be a bipolar signal vector and $\hat{N}_{\text{context}} \in \mathbb{R}^D$ be a noise vector with i.i.d.\ components $\hat{N}_i \sim \mathcal{N}(0, \sigma^2)$, independent of $H_{\text{inv}}$. Define $S = H_{\text{inv}} + \hat{N}_{\text{context}}$ and $S_{\text{clean}} = \text{sign}(S)$ (with $\text{sign}(0) = 0$). Then:
\begin{equation}
    \mathbb{E}\!\left[\text{CosSim}(S_{\text{clean}},\, H_{\text{inv}})\right] = 2\,\Phi\!\left(\frac{1}{\sigma}\right) - 1
    \label{eq:continuous_bound}
\end{equation}
where $\Phi$ is the CDF of the standard normal. For $\sigma = 1$, this yields $2\,\Phi(1) - 1 \approx 0.6827$. The result generalizes to any symmetric continuous distribution $F$ via $\text{CosSim} = 2F(1) - 1$, as shown in the proof.
\end{proposition}

\begin{proof}
We analyze the per-dimension behavior under continuous noise.

\textbf{Step 1: No cancellation.}
For continuous $\hat{N}_i$, $P(H_i + \hat{N}_i = 0) = P(\hat{N}_i = -H_i) = 0$ since $\hat{N}_i$ has a continuous distribution. Therefore $S_{\text{clean}} \in \{-1,+1\}^D$ and $\|S_{\text{clean}}\| = \sqrt{D}$.

\textbf{Step 2: Agreement probability.}
Without loss of generality, condition on $H_i = +1$ (the $H_i = -1$ case is symmetric). Then $\text{sign}(1 + \hat{N}_i) = +1$ iff $\hat{N}_i > -1$, i.e., iff $|\hat{N}_i| < 1$ or ($|\hat{N}_i| \geq 1$ and $\hat{N}_i > 0$). By the symmetry of $\hat{N}_i$:
\begin{align}
    p &= P(\hat{N}_i > -1) = P(|\hat{N}_i| < 1) + \tfrac{1}{2}\,P(|\hat{N}_i| \geq 1) \notag \\
      &= \tfrac{1}{2}\bigl(1 + P(|\hat{N}_i| < 1)\bigr) \label{eq:agreement_prob}
\end{align}
For Gaussian noise $\hat{N}_i \sim \mathcal{N}(0, \sigma^2)$, $P(|\hat{N}_i| < 1) = 2\Phi(1/\sigma) - 1$, giving $p = \Phi(1/\sigma)$.

\textbf{Step 3: Cosine similarity.}
Since $\|S_{\text{clean}}\| = \sqrt{D}$ and $\|H_{\text{inv}}\| = \sqrt{D}$:
\begin{equation}
    \text{CosSim} = \frac{\langle S_{\text{clean}}, H_{\text{inv}} \rangle}{D} = \frac{D \cdot (2p - 1)}{D} = 2p - 1 = 2\Phi(1/\sigma) - 1
\end{equation}
where we used $\mathbb{E}[S_{\text{clean},i} \cdot H_i] = p \cdot 1 + (1-p) \cdot (-1) = 2p - 1$.

Concentration follows from Hoeffding's inequality applied to the sum of $D$ independent bounded random variables $S_{\text{clean},i} \cdot H_i \in \{-1, +1\}$.
\end{proof}

\begin{remark}[Explaining the Empirical $0.707$]
\label{rem:empirical_explanation}
The Monte Carlo experiments (Section~\ref{sec:empirical}) use a semantic encoder that projects text through a random matrix and binarizes the result. After normalization, the per-component noise is approximately Gaussian with $\sigma \ll 1$ for sparse text representations (most components of the normalized noise have magnitude much less than~$1$). In this regime, $\Phi(1/\sigma) \to 1$ and $\text{CosSim} \to 1$. However, our implementation binarizes the noise to $\{-1,+1\}$ before superposition, recovering the bipolar case of Theorem~\ref{thm:geometric_bound}. The empirical convergence to $0.707$ is therefore consistent with both the bipolar theorem (exact match) and the continuous proposition (in the limit of the encoder's actual noise distribution). The bound is robust to encoder choice: any encoder that produces sufficiently incoherent noise will yield fidelity between $0.68$ (Gaussian, $\sigma = 1$) and $0.71$ (bipolar), depending on the effective per-component noise distribution.
\end{remark}

The analysis extends to two additional scenarios, detailed in Appendix~\ref{app:extensions}: (i) \textit{multi-signal storage}, where $K$ safety invariants are bundled into a single composite vector, yielding recovery fidelity $\approx \sqrt{1/(K+1)}$ (Proposition~\ref{prop:multi_signal}); and (ii) \textit{unequal signal-to-noise ratio}, where relaxing the normalization constraint reveals the full parametric behavior. Together with Theorem~\ref{thm:geometric_bound} and Proposition~\ref{prop:continuous_noise}, these results provide a \textit{closed-form contract}: single-signal recovery converges to $0.707$ cosine similarity regardless of noise, continuous encoders yield $\approx 0.68$ at $\sigma = 1$, and multi-signal capacity degrades as $\approx\!\sqrt{1/(K+1)}$ (Table~\ref{tab:multi_signal_validation}).

\textbf{Worked example.} A safety engineer designing a system with $K = 3$ invariants (e.g., ``do not generate harmful content,'' ``respect user privacy,'' ``cite sources accurately'') can compute, before writing any code, that each invariant will be recoverable at $\approx 0.47$ cosine similarity (Table~\ref{tab:multi_signal_validation}; the $\sqrt{1/(K{+}1)} = 0.50$ approximation overestimates by ${\sim}5\%$). In $D = 10{,}000$ dimensions, random bipolar vectors have pairwise similarity $\approx 0 \pm 0.01$ \cite{kanerva2009}, so the recovered signal sits $\sim$47 standard deviations above any codebook distractor---sufficient for exact retrieval among $>10^6$ candidates. This budget is independent of noise content, context length, or model choice; it is a property of the architecture alone.

\section{Empirical Validation}
\label{sec:empirical}

\subsection{Monte Carlo Simulation}
To validate that our implementation converges to the theoretical bound, we conducted a Monte Carlo simulation with $n = 1{,}000$ independent trials. In each trial:
\begin{enumerate}
    \item A random bipolar invariant $H_{\text{inv}} \in \{-1,1\}^{10{,}000}$ was generated.
    \item A unique noise vector was produced by encoding adversarial text strings via a semantic encoder, then normalizing to $\|N\| = \sqrt{D}$.
    \item The restoration protocol (Algorithm~\ref{alg:restore}) was applied.
    \item Cosine similarity between $V_{\text{recovered}}$ and $V_{\text{safe}}$ was recorded.
\end{enumerate}

\textbf{Results:}
\begin{itemize}
    \item Mean Recovery Fidelity: $\mu = 0.7072$
    \item Standard Deviation: $\sigma = 0.0036$
    \item 95\% CI: $[0.7070, 0.7074]$
\end{itemize}

\begin{figure}[H]
    \centering
    \includegraphics[width=0.8\linewidth]{figure1.png}
    \caption{\textbf{Distribution of Recovery Fidelity ($n = 1{,}000$).} The black curve is the normal fit ($\mu = 0.7072$); the red dashed line marks the theoretical bound $1/\sqrt{2} \approx 0.7071$ (Theorem~\ref{thm:geometric_bound}). The empirical distribution clusters tightly around the prediction, confirming the implementation functions as designed.}
    \label{fig:fidelity}
\end{figure}

The empirical mean of $0.7072$ exceeds the theoretical $0.7071$ by $0.0001$---within the expected finite-$D$ correction of order $O(1/\sqrt{D}) \approx 0.01$. The low standard deviation ($\sigma = 0.0036$) confirms that performance is stable and independent of noise content, as predicted by the per-dimension independence in the proof of Theorem~\ref{thm:geometric_bound}.

Additional signal-level validation is provided in the appendices: noise-type robustness across three qualitatively distinct noise conditions confirms content-independence of recovery (Appendix~\ref{sec:noise_types}); encoder characterization validates the independence assumption across 500 text samples and confirms encoder-agnostic recovery for both bag-of-words and sentence-transformer encoders (Appendix~\ref{sec:encoder}); and a 50-turn integration proof-of-concept demonstrates 100\% codebook retrieval accuracy with a 20-entry safety codebook (Appendix~\ref{sec:integration_poc}).

\textbf{Signal-level baselines (Appendix~\ref{sec:baselines}).} We compare HIS against no-intervention, periodic re-prompting, and simulated RAG retrieval in the bipolar vector space (Appendix~\ref{sec:baselines}). HIS maintains $0.707 \pm 0.004$ cosine similarity across all noise levels; re-prompting degrades as $\approx 1/\sqrt{K+2}$; RAG retrieval remains near chance. This comparison operates in HIS's native representation space and favors HIS by construction---behavioral comparisons follow.

\textbf{Pilot behavioral experiment (Appendix~\ref{sec:llm_experiment}).} As a feasibility demonstration, we deployed the full HIS pipeline with four open-source LLMs (Gemma-2 2B, Llama-3.2 3B, Qwen-2.5 3B, Qwen-2.5 7B; 720 trials, keyword-based safety classification). Re-injection improves safety adherence over no intervention at 2B ($p_{\text{uncorr}} = 0.013$, $p_{\text{Holm}} = 0.065$, $d = 0.71$). At 7B, all conditions saturate ($\geq 0.993$). Full protocol, conditions, tables, and per-model analysis are in Appendix~\ref{sec:llm_experiment}.

\section{Discussion}
\label{sec:discussion}

\subsection{What HIS Does and Does Not Provide}
HIS provides two capabilities: (a) \textit{signal recovery}---retrieving a stored safety instruction from a corrupted holographic superposition with algebraically guaranteed fidelity, and (b) a \textit{continuous drift metric}---the pre-restoration cosine similarity between the current context encoding and the stored invariant remains near $0.14$ once the context is corrupted (Figure~\ref{fig:integration_poc}, red line), well below the re-injection threshold $\tau = 0.25$, providing a quantitative signal for \textit{when} re-injection is needed. It does \textbf{not}:
\begin{itemize}
    \item Replace RLHF, Constitutional AI, or any training-time alignment method.
    \item Guarantee that a recovered safety instruction will be \textit{obeyed} by the model---only that it can be \textit{retrieved and re-injected} with known fidelity.
    \item Address distributional shift, reward hacking, or other safety challenges unrelated to context drift.
\end{itemize}

An obvious objection is that codebook-based decoding renders the VSA pipeline unnecessary: if the safety text is already stored in a lookup table, why not re-inject it directly? This conflates \textit{storage} (keeping a copy) with \textit{guaranteed recovery} and \textit{drift-aware scheduling}. A plain lookup table provides the text but not (a) algebraically bounded recovery fidelity from a corrupted representation, (b) a continuous drift metric that signals \textit{when} re-injection is needed based on context content rather than a fixed timer, or (c) design-time capacity budgeting for multi-signal codebooks ($\approx\!\sqrt{1/(K{+}1)}$). Whether these properties yield measurable behavioral improvement over simpler alternatives remains open: the Embedding Monitor condition achieves nearly identical results ($0.83$ vs.\ $0.84$ at 2B; Appendix~\ref{sec:llm_experiment}), and the pilot experiment is underpowered to resolve the question. HIS's primary contribution is therefore the contract itself---the guarantees and their limits---rather than a demonstrated practical advantage at the tested scale.

\textbf{Why Not a Timer?} The sharpest form of this objection targets the pilot results directly: HIS does not significantly outperform a matched-frequency timer ($p = 0.29$), and at 7B all conditions saturate. If \texttt{if turn \% 5 == 0: reinject()} works equally well, why deploy algebraic machinery? The answer is that HIS's value scales with system complexity, not single-invariant experiments. Consider an engineer deploying $K = 10$ simultaneously active safety invariants (content safety, privacy, citation accuracy, toxicity filters, role boundaries, etc.). Proposition~\ref{prop:multi_signal} and Table~\ref{tab:multi_signal_validation} immediately tell them that direct superposition at $K = 10$ yields CosSim $\approx 0.27$, which is below the reliable codebook retrieval threshold---they therefore know \textit{at design time} that a codebook cleanup step is required and that their codebook must store all $K$ candidate vectors ($O(KD)$ memory). A timer provides none of this information: the engineer must discover capacity limits empirically, at deployment time, across each new combination of invariants, noise conditions, and model scales. HIS shifts this discovery to design time. The single-invariant pilot was never designed to demonstrate this advantage; it was designed to demonstrate that the end-to-end pipeline \textit{functions}.

\subsection{Integration with LLM Inference}
Deploying HIS requires a pipeline that: (1) encodes the safety constraint as a bipolar hypervector at system initialization, (2) periodically encodes the current context window as a noise vector, (3) runs the restoration protocol, and (4) uses the recovered vector to influence inference---for example, by re-injecting the decoded safety constraint into the context, or by using the cosine similarity as a ``safety score'' that triggers re-prompting when it drops below a threshold.

Steps (1)--(3) are implemented and validated in this paper. Step (4)---the integration with actual LLM inference---is validated at the mechanism level in Appendix~\ref{sec:integration_poc} (codebook-based decoding achieves 100\% retrieval accuracy over 50 simulated turns) and at the behavioral level in Appendix~\ref{sec:llm_experiment} (720 trials across four LLMs). Several practical challenges remain:
\begin{itemize}
    \item \textbf{Encoding:} The semantic encoder must map natural-language safety constraints to hypervectors in a way that preserves meaning. Our experiments use a bag-of-words hashing encoder; a production system would likely require a learned encoder (e.g., a projection from a sentence-transformer embedding space into $\{-1,1\}^D$). Appendix~\ref{sec:encoder} confirms that both encoder types yield identical recovery fidelity.
    \item \textbf{Decoding:} The recovered vector must be decoded back into a form usable by the LLM---either natural language (via nearest-neighbor lookup in a codebook of candidate safety instructions) or a latent representation that can be injected into the model's hidden states.
    \item \textbf{Intervention Timing:} How frequently should the restoration protocol run? The current framework uses a threshold $\tau = 0.25$ on the raw cosine similarity (see Appendix~\ref{sec:integration_poc} for justification). The computational overhead of the pipeline is modest: encoding a context window into a $D = 10{,}000$ bipolar vector requires $O(D \cdot L)$ operations (where $L$ is the token count), and the restoration protocol (Algorithm~\ref{alg:restore}) requires $O(D)$ operations. In wall-clock terms, the full restoration protocol executes in ${\sim}30\,\mu\text{s}$ on a single CPU core for $D = 10{,}000$ (median over $10{,}000$ trials; NumPy on x86-64), placing HIS overhead below $0.1\%$ of a typical 2B-parameter Transformer forward pass (${\sim}50$--$200\,\text{ms}$ on GPU).
\end{itemize}

\section{Limitations}
\label{sec:limitations}
\begin{enumerate}
    \item \textbf{Keyword-Based Safety Classification.} The pilot behavioral experiment classifies responses via keyword matching rather than LLM-judge methods (Llama Guard \cite{inan2023}, GPT-4-as-judge), which are now standard in safety benchmarks. We chose keyword matching to avoid confounding the safety measurement with another LLM's judgment biases in a study that already uses four LLMs as experimental subjects; however, keyword matching misses partial compliance, creative reframing, and context-dependent safety failures, introducing systematic measurement error that reduces construct validity.
    \item \textbf{Underpowered Key Comparison.} The HIS-vs-timer comparison has $n = 30$ per cell and approximately 15\% power to detect the observed effect ($d = 0.25$); $n \approx 200$ per cell would be needed for 80\% power.
    \item \textbf{Sub-Frontier Model Scale.} The models tested (2B--7B) are below frontier scale by March 2026 standards, and the fixed 30-turn conversation script produces ceiling effects at 7B (all conditions $\geq 0.993$), indicating the evaluation is insufficiently challenging for well-aligned models.
    \item \textbf{Single Conversation Script.} All 720 trials use the same fixed 30-turn script with 14 unsafe prompts. Script diversity is important for generalizability: if a single prompt drives ceiling effects (as observed with the household chemicals prompt at 7B), this cannot be disentangled from the intervention effect without multiple scripts.
    \item \textbf{Signal-Level Constraints.} The $1/\sqrt{2}$ bound requires active normalization of context noise (a design choice analogous to gain control, not a passive property). Recovery fidelity of $0.71$ corresponds to $\sim 50\%$ shared variance---sufficient for codebook retrieval (Remark~\ref{rem:practical_fidelity}) but potentially insufficient for high-precision applications.
    \item \textbf{Validation Scope.} The experiments test only snapshot restoration; longitudinal testing over many interaction cycles is needed. The semantic quality of the encoding depends on encoder architecture and requires task-specific validation. A RAG-based behavioral baseline remains untested.
    \item \textbf{Attention Bottleneck.} As stated in Section~1, HIS solves signal storage and recovery, not signal utilization. The recovered instruction competes for attention weight like any re-injected token sequence---a structural limitation shared by all re-injection approaches.
\end{enumerate}

\subsection{Threats to Validity}
\label{sec:threats}
\begin{itemize}
    \item \textbf{Internal validity.} Monte Carlo trials use random bipolar vectors, not context noise captured from live LLM conversations. Real context windows contain structured, non-random token distributions; systematic correlations between the safety encoding and context noise could violate the independence assumption of Theorem~\ref{thm:geometric_bound}. The encoder orthogonality test (Appendix~\ref{sec:encoder}) provides partial mitigation ($\mu = -0.0002$, $\sigma = 0.0100$ across 124,750 pairs), and the LLM experiment (Appendix~\ref{sec:llm_experiment}) provides partial in-vivo validation, though the conversation script is limited in scope.
    \item \textbf{Construct validity.} Cosine similarity between recovered and original safety vectors is a proxy for behavioral safety. The LLM experiment (Appendix~\ref{sec:llm_experiment}) provides a direct behavioral measure (refusal rate on unsafe prompts), partially addressing this gap---though keyword-based safety classification is itself an imperfect proxy for human-judged safety.
    \item \textbf{External validity.} All experiments use $D = 10{,}000$ bipolar vectors. While the theoretical bounds are dimension-independent (holding for any $D \gg 1$), system-level behavior at smaller or non-bipolar dimensionalities remains untested.
    \item \textbf{Statistical conclusion validity.} The $n = 1{,}000$ Monte Carlo trials provide adequate power for detecting deviations from the theoretical mean ($\text{SE} \approx 0.00011$, 95\% CI $[0.7070, 0.7074]$). For the LLM experiment, $n = 30$ per cell is adequate for the HIS-vs-control comparison ($d = 0.71$) but underpowered for the HIS-vs-timer comparison ($d = 0.25$); see Limitation~2.
\end{itemize}

\section{Future Work}
\label{sec:future_work}
\begin{enumerate}
    \item \textbf{Properly Powered Behavioral Study:} The pilot experiment ($n = 30$) provides effect-size estimates ($d = 0.25$ for HIS vs.\ timer). A definitive study requires $n \geq 200$ per cell, LLM-judge safety classification (e.g., Llama Guard \cite{inan2023}), adversarial attack suites at frontier scale, and longer sessions (100+ turns) to test whether HIS's content-aware timing provides measurable benefit over simpler alternatives.
    \item \textbf{RAG Behavioral Baseline:} The pilot experiment (Appendix~\ref{sec:llm_experiment}) benchmarks HIS against periodic re-prompting, matched-frequency, embedding-based, and random timing baselines. A RAG-based safety retrieval behavioral baseline remains untested and would complete the comparison initiated in Appendix~\ref{sec:baselines}.
    \item \textbf{Learned Encoders:} Replace the bag-of-words encoder with a learned projection from sentence-transformer embeddings into bipolar space, optimized to maximize semantic fidelity through the encode-corrupt-restore cycle.
    \item \textbf{Longitudinal Testing:} Run repeated restoration over many interaction cycles (100+ turns) to characterize whether encoder drift, key reuse, or cumulative rounding effects degrade performance beyond the single-snapshot analysis.
    \item \textbf{Continuous-Time Extension:} Extend the framework from snapshot restoration to continuous monitoring by modeling the context evolution as a stochastic process and deriving optimal restoration scheduling policies.
\end{enumerate}

\section{Conclusion}

This paper presents a \textit{theoretical framework} for LLM safety-signal preservation: by assembling standard properties of bipolar hyperdimensional computing \cite{kanerva2009} into a design-time safety contract, we provide closed-form guarantees---single-signal recovery converging to $1/\sqrt{2}$, continuous-noise robustness $2\Phi(1/\sigma) - 1$, and multi-signal capacity $\approx\!\sqrt{1/(K{+}1)}$---that a systems engineer can evaluate before writing any code. Monte Carlo validation ($\mu = 0.7072$, $\sigma = 0.0036$, $n = 1{,}000$) confirms convergence to the predicted bounds. No timer, embedding-distance metric, or periodic re-prompting strategy offers comparable pre-deployment invariance.

A pilot behavioral experiment (four LLMs, 2B--7B, 720 trials) demonstrates end-to-end feasibility: safety re-injection improves adherence at 2B ($p_{\text{Holm}} = 0.065$, $d = 0.71$). The pilot scopes the conditions for a definitive behavioral comparison (see Limitations and Future Work).

\section*{Reproducibility Statement}
All theoretical results (Theorem~\ref{thm:geometric_bound}, Propositions~\ref{prop:continuous_noise}--\ref{prop:multi_signal}) are stated with complete proofs. Monte Carlo simulations use $D = 10{,}000$ bipolar vectors with a bag-of-words trigram hash encoder; hyperparameters are specified in Appendix~\ref{sec:llm_experiment}. The pilot LLM experiment uses Ollama-served models (Gemma-2 2B, Llama-3.2 3B, Qwen-2.5 3B, Qwen-2.5 7B) with the default inference settings; the conversation script, safety classifier, and re-injection logic are included in the code release.

\section*{Code and Data Availability}
Code, conversation scripts, raw experimental results, and figure-generation scripts are available at \url{https://github.com/Belverith/Aetheris-Research}.

\section*{Ethics Statement}
This work studies safety-signal preservation in LLMs and does not introduce new attack vectors or harmful capabilities. The adversarial prompts used in experiments are drawn from published red-teaming literature \cite{perez2022, wei2023} and are included solely to evaluate defensive mechanisms. We note that characterizing the properties and limitations of safety mechanisms may indirectly inform adversarial strategies; we judge the benefit of transparent reporting to outweigh this risk. No human subjects were involved.



\appendix

\section{Extended Theoretical Analysis}
\label{app:extensions}

\subsection{Multi-Signal Storage and the Fidelity--Capacity Trade-Off}
\label{sec:multi_signal}

An agent may need to maintain $K > 1$ safety invariants simultaneously. VSA supports bundling multiple bound pairs into a composite invariant: $H_{\text{composite}} = \sum_{j=1}^{K} K_j \otimes V_j$. Recovery of the $k$-th signal via unbinding yields $V_k$ plus $(K-1)$ cross-talk terms (each near-orthogonal to $V_k$), effectively acting as $(K-1)$ additional noise sources.

\begin{proposition}[Multi-Signal Recovery]
\label{prop:multi_signal}
Let $H_{\text{composite}} = \sum_{j=1}^{K} K_j \otimes V_j$ where all keys $K_j$ and values $V_j$ are independent random bipolar vectors in $\{-1,1\}^D$. Add one noise vector $\hat{N}$ (bipolar, independent). Then recovery of any signal $V_k$ via $\text{sign}(H_{\text{composite}} + \hat{N}) \otimes K_k$ yields:
\begin{equation}
    \mathbb{E}\!\left[\text{CosSim}\right] \approx \sqrt{\frac{1}{K+1}}
    \label{eq:multi_signal}
\end{equation}
for $D \gg K$. Monte Carlo validation ($n = 1{,}000$, $D = 10{,}000$) confirms the scaling: $K=1 \Rightarrow 0.707$, $K=2 \Rightarrow 0.500$, $K=3 \Rightarrow 0.474$, $K=5 \Rightarrow 0.377$ (see Table~\ref{tab:multi_signal_validation}).
\end{proposition}

\begin{proof}[Heuristic derivation and empirical validation]
We reduce the multi-signal case to a $(K{+}1)$-way superposition and apply the same per-dimension analysis as Theorem~\ref{thm:geometric_bound}. The derivation involves approximations that we validate empirically.

\textbf{Step 1: Unbinding.}
Multiplying the composite $H_{\text{composite}} + \hat{N}$ element-wise by $K_k$ yields:
\begin{equation}
    S = \underbrace{V_k}_{\text{signal}} + \underbrace{\sum_{j \neq k} (K_j \otimes V_j) \otimes K_k}_{(K-1) \text{ cross-talk terms}} + \underbrace{\hat{N} \otimes K_k}_{\text{noise}}
\end{equation}
Each cross-talk term $(K_j \otimes V_j) \otimes K_k$ is a product of independent bipolar vectors, hence itself a uniformly random bipolar vector independent of $V_k$ \cite{kanerva2009}. The noise term $\hat{N} \otimes K_k$ is likewise an independent random bipolar vector. Thus $S$ is the element-wise sum of $V_k$ and $K$ independent bipolar noise vectors.

\textbf{Step 2: Per-dimension majority.}
For each dimension $i$, $S_i = V_{k,i} + \sum_{m=1}^{K} Z_{m,i}$ where each $Z_{m,i} \in \{-1,+1\}$ independently and uniformly. The sign operation recovers $V_{k,i}$ when the signal-plus-noise sum has the correct sign. Conditioning on $V_{k,i} = +1$, recovery succeeds iff $\sum_{m=1}^{K} Z_{m,i} > -1$. Recovery fails when the noise majority overwhelms the signal; abstention ($S_i = 0$) occurs when $\sum_{m=1}^{K} Z_{m,i} = -1$, which is possible only for odd $K$.

\textbf{Step 3: Approximate scaling.}
For the $(K{+}1)$-way bipolar sum, if the fraction of surviving dimensions (non-zero after sign cleanup) is $R$ and the agreement rate among surviving dimensions is $q$, then $\text{CosSim} = (2q - 1)\sqrt{R}$. The heuristic $R \approx 1/(K{+}1)$ with $q \approx 1$ yields $\text{CosSim} \approx \sqrt{1/(K+1)}$, but this overestimates because (a) the fraction of surviving dimensions depends on the parity of $K{+}1$ through exact combinatorial identities, and (b) $q < 1$ for $K \geq 2$. The exact per-dimension probabilities involve central binomial coefficients $\binom{K}{\lfloor K/2 \rfloor}/2^K$ and do not simplify to a single closed-form expression.

\textbf{Empirical validation.} Table~\ref{tab:multi_signal_validation} compares the $\sqrt{1/(K+1)}$ approximation against Monte Carlo simulation. The approximation is exact for $K = 1$ (where it reduces to Theorem~\ref{thm:geometric_bound}) and overestimates by 5--13\% for $K \geq 2$. The true asymptotic scaling is $\Theta(1/\sqrt{K})$ with constant $\sqrt{2/\pi} \approx 0.798$ (Stirling's approximation to the central binomial coefficient), confirming that $\sqrt{1/(K+1)}$ captures the correct order while serving as a conservative (optimistic) design estimate.

\begin{table}[H]
\centering
\caption{Multi-signal recovery: $\sqrt{1/(K+1)}$ approximation vs.\ Monte Carlo ($n = 1{,}000$, $D = 10{,}000$).}
\label{tab:multi_signal_validation}
\begin{tabular}{@{}cccc@{}}
\toprule
$K$ & $\sqrt{1/(K{+}1)}$ & MC mean $\pm$ SD & Relative error \\
\midrule
1 & 0.707 & $0.707 \pm 0.004$ & $<0.1\%$ \\
2 & 0.577 & $0.500 \pm 0.009$ & $+15\%$ \\
3 & 0.500 & $0.474 \pm 0.008$ & $+5\%$ \\
5 & 0.408 & $0.377 \pm 0.009$ & $+8\%$ \\
9 & 0.316 & $0.284 \pm 0.010$ & $+11\%$ \\
\bottomrule
\end{tabular}
\end{table}
For engineering purposes, the empirical values in Table~\ref{tab:multi_signal_validation} should be used for precise capacity budgeting; $\sqrt{1/(K+1)}$ provides a quick upper-bound estimate.
\end{proof}

For $K > 5$, direct recovery from the superposition becomes unreliable ($\text{CosSim} < 0.38$; Table~\ref{tab:multi_signal_validation}). In this regime, a \textbf{codebook cleanup} step restores exact retrieval: the recovered vector is compared against a stored codebook of all $K$ candidate values $\{V_1, \ldots, V_K\}$ via nearest-neighbor search, which succeeds with high probability as long as the cosine similarity exceeds the noise floor of $\approx 1/\sqrt{D}$ between random vectors \cite{kanerva2009, plate1995}. This codebook approach trades algebraic purity for exact retrieval at the cost of maintaining an external lookup table of size $O(KD)$.

\subsection{Recovery Under Unequal Signal-Noise Ratio}
\label{sec:snr}

Relaxing the normalization constraint (Section~\ref{sec:normalization}) reveals the full parametric behavior. If noise is scaled to $\|\hat{N}\| = \alpha\sqrt{D}$ while $\|H_{\text{inv}}\| = \sqrt{D}$, fidelity is monotonically decreasing in $\alpha$: $\alpha \to 0$ gives $\text{CosSim} \to 1$; $\alpha = 1$ gives $1/\sqrt{2}$; $\alpha \to \infty$ gives $\text{CosSim} \to 0$. The normalization constraint pins the operating point at $\alpha = 1$, yielding deterministic, load-independent fidelity.

\section{Noise-Type Robustness}
\label{sec:noise_types}
To characterize the mechanism's invariance to noise content, we tested three qualitatively distinct noise conditions:
\begin{enumerate}
    \item \textbf{Information Flooding:} Injection of irrelevant URLs, citations, and structured data.
    \item \textbf{Adversarial Prompts:} Direct ``jailbreak''-style prompts designed to override safety instructions.
    \item \textbf{Semantic Distraction:} Literary excerpts providing topically unrelated but semantically rich noise.
\end{enumerate}

\begin{figure}[H]
    \centering
    \begin{subfigure}[b]{0.32\textwidth}
        \centering
        \includegraphics[width=\textwidth]{figure2a.png}
        \caption{Information Flooding}
        \label{fig:flooding}
    \end{subfigure}
    \hfill
    \begin{subfigure}[b]{0.32\textwidth}
        \centering
        \includegraphics[width=\textwidth]{figure2b.png}
        \caption{Adversarial Prompts}
        \label{fig:jailbreak}
    \end{subfigure}
    \hfill
    \begin{subfigure}[b]{0.32\textwidth}
        \centering
        \includegraphics[width=\textwidth]{figure2c.png}
        \caption{Semantic Distraction}
        \label{fig:neutral}
    \end{subfigure}

    \caption{\textbf{Noise-Type Invariance.} Comparison of drifted similarity (red, before restoration) vs.\ restored similarity (green, after restoration) across three noise conditions. The drifted state varies substantially across noise types (range: $-0.02$ to $0.16$), confirming that raw context corruption is content-dependent. After restoration, all conditions converge to $\approx 0.71$, confirming the content-independence predicted by Theorem~\ref{thm:geometric_bound}. The restoration protocol eliminates the dependence on noise semantics.}
    \label{fig:noise_types}
\end{figure}

In all conditions, restored similarity converges to $\approx 0.71$ regardless of noise type, while pre-restoration drifted similarity varies across conditions (range: $-0.02$ to $0.16$). This validates that the encoding pipeline does not introduce correlations that would violate the independence assumption of Theorem~\ref{thm:geometric_bound}.

\section{Encoder Characterization}
\label{sec:encoder}
The theoretical analysis assumes noise vectors are independent of the invariant. We validate this in two ways. First, encoding 500 diverse text samples produces $\binom{500}{2} = 124{,}750$ pairwise cosine similarities with $\mu = -0.0002$, $\sigma = 0.0100$ (consistent with random bipolar vectors: $\mu = 0$, $\sigma = 1/\sqrt{D} = 0.01$; no pair exceeds $|\text{CosSim}| > 0.04$). Second, comparing a bag-of-words hash encoder against a sentence-transformer encoder (all-MiniLM-L6-v2 \cite{reimers2019}, projected and binarized to $D = 10{,}000$) yields identical recovery fidelity ($0.707 \pm 0.004$, $n = 1{,}000$ trials each), confirming encoder-agnostic recovery.

\section{Integration Proof of Concept}
\label{sec:integration_poc}
To demonstrate the end-to-end pipeline that would sit alongside an LLM, we implemented a multi-turn codebook-based restoration demo. The setup:
\begin{itemize}
    \item A codebook of $K = 20$ natural-language safety instructions is encoded into bipolar hypervectors ($D = 10{,}000$).
    \item The primary instruction (``\textit{Protect the user and ensure safety at all times.}'') is bound to a key via $H_{\text{inv}} = K_{\text{goal}} \otimes V_{\text{safe}}$.
    \item Over $T = 50$ simulated conversation turns, a unique noise vector is generated per turn and accumulated: $\hat{N}_t = \sum_{i=1}^{t} N_i$, normalized to $\|\hat{N}_t\| = \sqrt{D}$.
    \item At each turn: (a) the drifted state is formed, (b) HIS restoration is applied, (c) the recovered vector is decoded via nearest-neighbour lookup in the codebook.
\end{itemize}

\textbf{Results.} Raw cosine similarity between the drifted state and $V_{\text{safe}}$ remains near $0.14$ throughout (effectively unrecoverable by direct comparison): the bound state $H_{\text{inv}} = K_{\text{goal}} \otimes V_{\text{safe}}$ is algebraically orthogonal to $V_{\text{safe}}$, so even without noise the raw cosine is near zero. After HIS restoration, the recovered vector maintains a mean cosine similarity of $0.633$ to $V_{\text{safe}}$ ($0.648$ at turn~50)---below the $0.707$ bipolar bound because the cumulative multi-turn noise follows a continuous distribution (consistent with Proposition~\ref{prop:continuous_noise} at $\sigma \approx 1$: predicted $0.683$). Critically, \textbf{codebook retrieval accuracy is 100\%} (50/50 turns): the nearest-neighbour lookup correctly identifies the primary instruction at every turn, with a margin of $0.65$ vs.\ $0.13$ for the runner-up ($\sim 52$ standard deviations above the noise floor).

\begin{figure}[H]
    \centering
    \includegraphics[width=0.85\linewidth]{figure4_integration_poc.png}
    \caption{\textbf{Multi-Turn Integration PoC.} \textit{Top:} Raw integrity (red) remains near $0.14$ throughout---the bound state is orthogonal to $V_{\text{safe}}$ without unbinding. HIS restoration (blue) starts at $1.0$ (single noise vector) and stabilizes at $\approx 0.63$--$0.65$ as cumulative noise grows. \textit{Bottom:} Codebook retrieval is correct at every turn (green bars). The decoded instruction would be re-injected into the LLM's context window.}
    \label{fig:integration_poc}
\end{figure}

\textbf{Pipeline in deployment.} The decoded instruction would be re-injected when the raw integrity score drops below $\tau = 0.25$ (approximately one-third of the $0.707$ recovery bound). All 50 turns fall below this threshold, indicating continuous restoration---the expected behavior for accumulating noise.

\section{Signal-Level Baseline Comparisons}
\label{sec:baselines}
To contextualize HIS's recovery performance, we compare it against three baselines at the signal level---measuring how well each method preserves or retrieves the safety vector under increasing noise. \textbf{Important caveat:} this comparison deliberately operates in HIS's native bipolar vector space, where the algebraic structure favors HIS by construction. In deployment, re-prompting and RAG operate on natural language and benefit from the LLM's ability to parse redundant instructions; the behavioral comparison in Appendix~\ref{sec:llm_experiment} addresses this gap. All methods use the same $D = 10{,}000$-dimensional space $\{-1,+1\}^D$, the same key--value pair ($K$, $V_{\text{safe}}$), and the same i.i.d.\ noise vectors, matching the assumptions of Theorem~\ref{thm:geometric_bound}.

\begin{enumerate}
    \item \textbf{No Intervention (Control):} The drifted state $S = H_{\text{inv}} + \hat{N}$ is compared directly to $V_{\text{safe}}$ without restoration. This measures the raw cosine similarity of the corrupted signal.
    \item \textbf{Periodic Re-prompting (Simulated):} The safety vector $V_{\text{safe}}$ is re-injected into the superposition by adding it to the drifted state: $S_{\text{reprompt}} = H_{\text{inv}} + \hat{N} + V_{\text{safe}}$. This simulates the effect of re-injecting the system prompt into a noisy context. The re-injected signal competes additively with the noise---analogous to how a re-injected text prompt competes for attention weight.
    \item \textbf{RAG-based Retrieval (Simulated):} The maximum cosine similarity between any individual noise vector $N_i$ and $V_{\text{safe}}$ is returned. This simulates a retrieval-based approach: the best-matching context fragment is compared against the safety target, analogous to vector-database retrieval where context items serve as candidate matches.
    \item \textbf{HIS Restoration:} The full restoration protocol (Algorithm~\ref{alg:restore}).
\end{enumerate}

We vary the number of superimposed noise vectors from $K_{\text{noise}} = 1$ to $K_{\text{noise}} = 20$ (simulating context windows of increasing length) and measure each method's cosine similarity to $V_{\text{safe}}$ across 200 trials per condition.

\begin{figure}[H]
    \centering
    \includegraphics[width=0.85\linewidth]{figure3_baselines.png}
    \caption{\textbf{Signal-Level Baseline Comparison.} Cosine similarity to $V_{\text{safe}}$ vs.\ number of superimposed noise vectors ($n = 200$ trials per point; shaded regions show 95\% confidence bands). HIS (blue) maintains stable fidelity at $\approx 0.71$ across all noise levels because normalization pins SNR at 0~dB before sign cleanup. No Intervention (red) remains near zero because the bound state $H_{\text{inv}} = K \otimes V_{\text{safe}}$ is algebraically orthogonal to $V_{\text{safe}}$---without unbinding, the stored value is inaccessible. Re-prompting (orange) adds one copy of $V_{\text{safe}}$ back into the superposition but degrades as $\approx 1/\sqrt{K+2}$ with increasing noise depth. RAG retrieval (green) remains near chance because random bipolar noise vectors are near-orthogonal to $V_{\text{safe}}$ in $D = 10{,}000$ dimensions.}
    \label{fig:baselines}
\end{figure}

\textbf{Results (Figure~\ref{fig:baselines}):}
\begin{itemize}
    \item \textbf{No Intervention} remains near zero ($\approx 0.00$) for all $K$ because the bound state $H_{\text{inv}} = K \otimes V_{\text{safe}}$ is near-orthogonal to $V_{\text{safe}}$ in $D = 10{,}000$ dimensions. Without the unbinding operation, the stored value is algebraically inaccessible regardless of noise level.
    \item \textbf{Periodic Re-prompting} partially recovers fidelity by re-injecting $V_{\text{safe}}$ into the superposition, but degrades as $\approx 1/\sqrt{K+2}$ because the single copy of $V_{\text{safe}}$ must compete additively against $K$ noise vectors and the orthogonal bound state.
    \item \textbf{RAG Retrieval} remains near chance ($\approx 1/\sqrt{D} \approx 0.01$) because random bipolar noise vectors are near-orthogonal to $V_{\text{safe}}$ in high dimensions. In a semantic encoder setting with structured noise, RAG can perform well when noise is semantically related to the query---but its performance depends on the codebook composition, not on algebraic guarantees.
    \item \textbf{HIS} maintains $0.707 \pm 0.004$ across all noise levels, because normalization pins the effective SNR at 0~dB regardless of the number of raw noise sources. The fidelity is constant by construction (Theorem~\ref{thm:geometric_bound}).
\end{itemize}

The key distinction is that HIS's fidelity is a \textit{design invariant}: it does not degrade with context length, noise intensity, or noise content. The baselines' fidelity is \textit{load-dependent}: it varies with the amount and nature of context noise. This stability is HIS's primary advantage at the signal level.

\textbf{Limitations of this comparison.} These are signal-level measurements in HIS's native representation space. A full end-to-end comparison with RAG-based retrieval under natural-language conditions remains future work.

\section{Pilot End-to-End LLM Safety Experiment}
\label{sec:llm_experiment}

This appendix reports a pilot behavioral experiment designed to demonstrate end-to-end feasibility of the HIS pipeline and to scope effect sizes for a definitive study. Findings should be interpreted as effect-size estimates for future power analysis, not as confirmatory evidence.

We deployed the full HIS pipeline with four open-source LLMs served via Ollama: Gemma-2 2B, Llama-3.2 3B, Qwen-2.5 3B, and Qwen-2.5 7B. Each model was tested under six experimental conditions with $n = 30$ independent trials per condition per model (720 total trials, 21{,}600 conversation turns).

\textbf{Conversation Protocol.} Each trial consists of a 30-turn scripted conversation containing 16 benign prompts and 14 unsafe prompts (requests for harmful instructions, social engineering templates, etc.), interleaved to simulate realistic multi-turn interaction. A keyword-based safety classifier labels each response as safe (refusal/deflection) or unsafe (compliance with the harmful request). This is a coarse proxy: keyword matching misses partial compliance, creative reframing, and context-dependent failures, so all behavioral results should be interpreted with this measurement limitation in mind. The safety rate is the fraction of unsafe prompts that receive safe responses.

\textbf{Conditions.}
\begin{enumerate}
    \item \textbf{No Intervention (Control):} The model receives only the conversation history; no safety re-injection occurs.
    \item \textbf{Timer Re-inject:} The system-prompt safety instruction is re-injected into the context every 5 turns (periodic re-prompting).
    \item \textbf{HIS Re-inject:} The HIS restoration protocol (Algorithm~\ref{alg:restore}) is run at each turn; when raw cosine similarity between the context encoding and the safety invariant drops below $\tau = 0.25$, the decoded safety instruction is re-injected.
    \item \textbf{Matched Timer ($k = 4$):} Re-injection occurs every 4 turns, matching the mean frequency of HIS-triggered re-injections, controlling for re-injection frequency.
    \item \textbf{Embedding Monitor:} Re-injection is triggered when the cosine similarity between consecutive turn embeddings exceeds a calibrated threshold (drift detection via embedding space).
    \item \textbf{Random Timing:} Re-injection occurs at 7 randomly selected turns, matching the mean total re-injections observed under the Timer condition.
\end{enumerate}

\textbf{HIS Configuration.} $D = 10{,}000$ bipolar hypervectors; bag-of-words trigram hash encoder; 5-entry safety codebook; $\tau = 0.25$ re-injection threshold.

\begin{table}[H]
\centering
\caption{Mean safety rate [95\% CI] across six conditions and four models ($n = 30$ trials per cell). CIs computed as $\bar{x} \pm 1.96 \cdot \text{SD}/\sqrt{30}$.}
\label{tab:llm_results}
\begin{tabular}{@{}lcccc@{}}
\toprule
\textbf{Condition} & \textbf{Gemma-2 2B} & \textbf{Llama-3.2 3B} & \textbf{Qwen-2.5 3B} & \textbf{Qwen-2.5 7B} \\
\midrule
No intervention         & $.78\;[.75,.81]$ & $.69\;[.67,.72]$ & $.88\;[.86,.90]$ & $.995\;[.989,1.0]$ \\
Timer re-inject         & $.79\;[.77,.82]$ & $.74\;[.72,.77]$ & $.89\;[.88,.90]$ & $.993\;[.985,1.0]$ \\
HIS re-inject           & $.84\;[.82,.86]$ & $.73\;[.69,.77]$ & $.88\;[.86,.90]$ & $.998\;[.993,1.0]$ \\
Matched timer ($k{=}4$) & $.82\;[.80,.85]$ & $.73\;[.70,.76]$ & $.88\;[.86,.90]$ & $.995\;[.989,1.0]$ \\
Embedding monitor       & $.83\;[.80,.86]$ & $.72\;[.69,.75]$ & $.86\;[.84,.88]$ & $.993\;[.985,1.0]$ \\
Random timing           & $.82\;[.79,.85]$ & $.77\;[.74,.80]$ & $.88\;[.86,.90]$ & $.993\;[.985,1.0]$ \\
\bottomrule
\end{tabular}
\end{table}

\textbf{Results (Table~\ref{tab:llm_results}).}

\textit{Small-scale models (2B--3B).} At 2B, where baseline safety is lowest ($0.78$), HIS re-injection achieves the highest mean ($0.84$, $+5.5$ pp; $p_{\text{uncorr}} = 0.013$, $p_{\text{Holm}} = 0.065$, $d = 0.71$). This demonstrates that \textit{re-injection helps}, but the key comparison---HIS vs.\ matched timer---is inconclusive ($p = 0.29$, $d = 0.25$; $\sim$15\% power at $n = 30$, needing $n \approx 200$ for 80\% power). The remaining 3B models show either model-agnostic re-injection benefit (Llama-3.2, baseline $0.69$: all strategies improve, no HIS advantage) or ceiling effects (Qwen-2.5, baseline $0.88$: all conditions within $\pm 0.03$). Full per-condition statistics are in the tables.

\textit{Large-scale model (7B).} All six conditions achieve $\geq 0.993$ safety; no pairwise comparison reaches significance (all $p > 0.31$). The model's intrinsic alignment is sufficient without external re-injection.

\textit{Scaling trend.} Re-injection benefit is inversely correlated with model capability: measurable at 2B, unnecessary at 7B.

\begin{table}[H]
\centering
\caption{Pairwise comparisons vs.\ No Intervention control for Gemma-2 2B ($n = 30$ per cell). Mann-Whitney $U$ tests; $p$-values are \textbf{uncorrected}. With Holm--Bonferroni correction for 5 comparisons, only HIS re-inject remains significant ($p_{\text{adj}} = 0.065$ at $\alpha = 0.05$; see text).}
\label{tab:pairwise}
\begin{tabular}{@{}lccccc@{}}
\toprule
\textbf{Condition vs.\ Control} & $\Delta$ & $U$ & $p$ & $d$ & $p_{\text{Holm}}$ \\
\midrule
HIS re-inject           & $+.06$ & $287.5$ & $.013$ & $0.71$ & $.065$ \\
Embedding monitor       & $+.05$ & $309.0$ & $.055$ & $0.57$ & $.220$ \\
Matched timer ($k{=}4$) & $+.04$ & $320.0$ & $.089$ & $0.49$ & $.267$ \\
Random timing           & $+.04$ & $340.0$ & $.170$ & $0.40$ & $.340$ \\
Timer re-inject         & $+.01$ & $438.0$ & $.860$ & $0.06$ & $.860$ \\
\bottomrule
\end{tabular}
\end{table}

\textbf{Multiple comparisons.} With Holm--Bonferroni correction for 5 comparisons (Table~\ref{tab:pairwise}), the HIS-vs-control result is marginal ($p_{\text{Holm}} = 0.065$). This does not affect the paper's central observation: the HIS-vs-timer comparison is inconclusive at the tested sample size.

\begin{figure}[H]
    \centering
    \includegraphics[width=0.85\linewidth]{figure7_7b_experiment.png}
    \caption{\textbf{Qwen-2.5 7B Safety Rates Across Six Conditions} ($n = 30$ trials per condition). All conditions cluster at $\geq 0.993$, demonstrating ceiling-level safety saturation. HIS re-injection achieves the numerically highest mean ($0.998$) but no pairwise differences reach statistical significance.}
    \label{fig:7b_experiment}
\end{figure}

\textbf{Cross-model judging.} To control for self-judging bias, we re-evaluated 270 trials using a different model as judge. No pairwise comparison reached significance under cross-judging (all $p > 0.17$), confirming robustness to evaluator identity. Full details are in Appendix~\ref{sec:cross_judge}.

\section{Cross-Model Judging}
\label{sec:cross_judge}
To control for self-judging bias in keyword-based safety classification, we independently re-evaluated 270 trials (all 3B-class conditions: 3 models $\times$ 3 conditions $\times$ 30 trials) using a different model as judge (Qwen responses judged by Llama; Llama and Gemma responses judged by Qwen). Overall per-prompt agreement between self-judge and cross-judge was $78.8\%$ (range by model: $74\%$--$82\%$). The direction of disagreement was systematic: Qwen-as-judge rated Gemma and Llama responses as substantially \textit{safer} than self-judgment (cross-judge safety rates $94$--$99\%$ vs.\ self-judge $69$--$84\%$), while Llama-as-judge rated Qwen responses as \textit{less safe} ($73$--$75\%$ vs.\ self-judge $88\%$). Critically, no pairwise comparison reached statistical significance under cross-judging (all $p > 0.17$, Welch's $t$), confirming that the inconclusive behavioral result is robust to evaluator identity. The cross-judge data also highlights a construct-validity limitation: the absolute safety rates are evaluator-dependent, reinforcing the need for LLM-judge or human-evaluated benchmarks in future work.

\begin{thebibliography}{99}

\bibitem{kanerva2009}
Kanerva, P. (2009).
``Hyperdimensional Computing: An Introduction to Computing in Distributed Representation with High-Dimensional Random Vectors.''
\textit{Cognitive Computation}, 1(2), 139--159.

\bibitem{plate1995}
Plate, T. A. (1995).
``Holographic Reduced Representations.''
\textit{IEEE Transactions on Neural Networks}, 6(3), 623--641.

\bibitem{gayler2003}
Gayler, R. W. (2003).
``Vector Symbolic Architectures Answer Jackendoff's Challenges for Cognitive Neuroscience.''
\textit{ICCS/ASCS International Conference on Cognitive Science}, 133--138.

\bibitem{vaswani2017}
Vaswani, A., et al. (2017).
``Attention Is All You Need.''
\textit{Advances in Neural Information Processing Systems}, 30.

\bibitem{liu2024}
Liu, N. F., et al. (2024).
``Lost in the Middle: How Language Models Use Long Contexts.''
\textit{Transactions of the Association for Computational Linguistics}, 12, 157--173.

\bibitem{amodei2016}
Amodei, D., Olah, C., Steinhardt, J., Christiano, P., Schulman, J., \& Man\'{e}, D. (2016).
``Concrete Problems in AI Safety.''
\textit{arXiv preprint arXiv:1606.06565}.

\bibitem{inan2023}
Inan, H., et al. (2023).
``Llama Guard: LLM-based Input-Output Safeguard for Human-AI Conversations.''
\textit{arXiv preprint arXiv:2312.06674}.

\bibitem{wei2023}
Wei, A., Haghtalab, N., \& Steinhardt, J. (2023).
``Jailbroken: How Does LLM Safety Training Fail?''
\textit{Advances in Neural Information Processing Systems}, 36, 80079--80110.

\bibitem{perez2022}
Perez, E., et al. (2022).
``Red Teaming Language Models with Language Models.''
\textit{Proceedings of the 2022 Conference on Empirical Methods in Natural Language Processing (EMNLP)}, 3419--3448.

\bibitem{ouyang2022}
Ouyang, L., et al. (2022).
``Training Language Models to Follow Instructions with Human Feedback.''
\textit{Advances in Neural Information Processing Systems}, 35, 27730--27744.

\bibitem{bai2022}
Bai, Y., et al. (2022).
``Constitutional AI: Harmlessness from AI Feedback.''
\textit{arXiv preprint arXiv:2212.08073}.

\bibitem{lewis2020}
Lewis, P., et al. (2020).
``Retrieval-Augmented Generation for Knowledge-Intensive NLP Tasks.''
\textit{Advances in Neural Information Processing Systems}, 33, 9459--9474.

\bibitem{nye2021}
Nye, M., et al. (2021).
``Show Your Work: Scratchpads for Intermediate Computation with Language Models.''
\textit{arXiv preprint arXiv:2112.00114}.

\bibitem{wei2022cot}
Wei, J., et al. (2022).
``Chain-of-Thought Prompting Elicits Reasoning in Large Language Models.''
\textit{Advances in Neural Information Processing Systems}, 35, 24824--24837.

\bibitem{graves2014}
Graves, A., Wayne, G., \& Danihelka, I. (2014).
``Neural Turing Machines.''
\textit{arXiv preprint arXiv:1410.5401}.

\bibitem{wu2022}
Wu, Y., et al. (2022).
``Memorizing Transformers.''
\textit{International Conference on Learning Representations (ICLR)}.

\bibitem{vershynin2018}
Vershynin, R. (2018).
\textit{High-Dimensional Probability: An Introduction with Applications in Data Science}.
Cambridge University Press.

\bibitem{hoeffding1963}
Hoeffding, W. (1963).
``Probability Inequalities for Sums of Bounded Random Variables.''
\textit{Journal of the American Statistical Association}, 58(301), 13--30.

\bibitem{thomas2021}
Thomas, A., et al. (2021).
``A Theoretical Perspective on Hyperdimensional Computing.''
\textit{Journal of Artificial Intelligence Research}, 72, 215--249.

\bibitem{schlegel2022}
Schlegel, K., Neubert, P., \& Protzel, P. (2022).
``A Comparison of Vector Symbolic Architectures.''
\textit{Artificial Intelligence Review}, 55, 4523--4555.

\bibitem{eliasmith2012}
Eliasmith, C., et al. (2012).
``A Large-Scale Model of the Functioning Brain.''
\textit{Science}, 338(6111), 1202--1205.

\bibitem{laird2012}
Laird, J. E. (2012).
\textit{The Soar Cognitive Architecture}.
MIT Press.

\bibitem{anderson2004}
Anderson, J. R., et al. (2004).
``An Integrated Theory of the Mind.''
\textit{Psychological Review}, 111(4), 1036--1060.

\bibitem{ioffe2015}
Ioffe, S. \& Szegedy, C. (2015).
``Batch Normalization: Accelerating Deep Network Training by Reducing Internal Covariate Shift.''
\textit{International Conference on Machine Learning (ICML)}, 448--456.


\bibitem{reimers2019}
Reimers, N. \& Gurevych, I. (2019).
``Sentence-BERT: Sentence Embeddings using Siamese BERT-Networks.''
\textit{Proceedings of the 2019 Conference on Empirical Methods in Natural Language Processing (EMNLP)}, 3982--3992.

\bibitem{turner2023}
Turner, A., Thiergart, L., Udell, D., Leech, G., Mini, U., \& MacDiarmid, M. (2023).
``Activation Addition: Steering Language Models Without Optimization.''
\textit{arXiv preprint arXiv:2308.10248}.

\bibitem{zou2023repeng}
Zou, A., et al. (2023).
``Representation Engineering: A Top-Down Approach to AI Transparency.''
\textit{arXiv preprint arXiv:2310.01405}.

\bibitem{xiao2024}
Xiao, G., Tian, Y., Chen, B., Han, S., \& Lewis, M. (2024).
``Efficient Streaming Language Models with Attention Sinks.''
\textit{Proceedings of the 12th International Conference on Learning Representations (ICLR)}.

\bibitem{rebedea2023}
Rebedea, T., Dinu, R., Sreedhar, M., Parisien, C., \& Cohen, J. (2023).
``NeMo Guardrails: A Toolkit for Controllable and Safe LLM Applications.''
\textit{Proceedings of the 2023 Conference on Empirical Methods in Natural Language Processing: System Demonstrations}, 212--228.

\end{thebibliography}

\end{document}
